\documentclass[11pt,a4paper]{article}
\usepackage[margin=24mm]{geometry}
\usepackage{amsmath,amssymb}
\title{Class Field Theory VIII.6.13}
\author{J. S. Milne}
\date{Fixed author PDF v4.03}
\begin{document}
\maketitle
\section*{Original statement}
\noindent\emph{Source \texttt{CFT.pdf}, physical page 265.}\par\medskip
% BEGIN CFT_VIII613_STATEMENT
\noindent\textsc{Lemma 6.13}
\emph{Let $T$ be finite set of real or finite primes of $K$, and let
$b\in K^\times$. If $T$ has an even number of elements and $b$ does not
become a square in $K_v^\times$ for any $v\in T$, then there exists an
$a\in K^\times$ such that}
\[
(a,b)_v =
\begin{cases}
-1 & \text{for }v\in T\\
 1 & \text{otherwise.}
\end{cases}
\]
% END CFT_VIII613_STATEMENT

\section*{Original proof}
\noindent\emph{Source \texttt{CFT.pdf}, physical pages 265--266.}\par\medskip
% BEGIN CFT_VIII613_PROOF
\noindent\textsc{Proof.} (Following Tate, 1976, 5.2). Let $L$ be the
composite of all abelian extensions of $K$ of exponent 2, and let
$G=\operatorname{Gal}(L/K)$. By class field theory,
\[
G\simeq \boldsymbol{C}_K/2\boldsymbol{C}_K
 \simeq \mathbb{I}/K^\times\cdot\mathbb{I}^2.
\]
The cohomology sequence of
\[
0\longrightarrow\mu_2\longrightarrow L^\times
 \xrightarrow{x\mapsto x^2}L^{\times2}\longrightarrow0
\]
is an exact sequence
\[
K^\times\xrightarrow{x\mapsto x^2}K^\times\cap L^{\times2}
 \longrightarrow\operatorname{Hom}_{\mathrm{conts}}(G,\mu_2)
 \longrightarrow0.
\]
Every element of $K^\times$ becomes a square in $L$, and so we have an
isomorphism
\[
K^\times/K^{\times2}\longrightarrow
\operatorname{Hom}_{\mathrm{conts}}(\mathbb{I}/K^\times\cdot\mathbb{I}^2,\mu_2).
\]
This map sends $a\in K^\times$ to the continuous homomorphism
\[
(c_v)\longmapsto\prod(a,c_v)_v
\]
(because of the relation between the Hilbert symbol and the local Artin map).
Thus, finding $a$ is equivalent to finding a homomorphism
$f:\mathbb{I}/\mathbb{I}^2\longrightarrow\mu_2$ such that
\begin{enumerate}
\renewcommand{\labelenumi}{(\alph{enumi})}
\item $f=1$ on principal id\`eles;
\item
\[
f(1,\ldots,1,i_v(b),1,\ldots,1)=
\begin{cases}
-1 & \text{for }v\in T\\
 1 & \text{otherwise.}
\end{cases}
\]
\end{enumerate}
where $i_v$ is the inclusion $K\hookrightarrow K_v$. For each $v$,
let $B_v$ be the $\mathbb{F}_2$-subspace of $K_v^\times/K_v^{\times2}$
generated by $i_v(b)$, and let
$B=\prod B_v\subset\mathbb{I}/\mathbb{I}^2$.
Because $i_v(b)$ is not a square for $v\in T$, there exists a linear
form (automatically continuous) $f_1:B\longrightarrow\mu_2$ satisfying
condition (b), and $f_1$ will extend to a continuous linear form on
$\mathbb{I}/\mathbb{I}^2$ satisfying (a) if and only if $f_1$ takes the
value $1$ on every principal id\'ele in $B$. The value of $f_1$ on the
principal id\'ele of $b$ is $\prod_{v\in T}-1$, which is $1$ because of
our assumption that $T$ contains an even number of elements. Let
$c\in K^\times$ be such that its id\'ele lies in $B$. Then, for every $v$,
$i_v(c)=1$ or $i_v(b)$ in $K_v^\times/K_v^{\times2}$. Therefore,
$i_v(c)$ becomes a square in $K_v[\sqrt{b}]$ for all $v$, which (by 1.1)
implies that $c$ is a square in $K[\sqrt{b}]$. Hence $c=1$ or $b$ in
$K^\times/K^{\times2}$, and so $f_1$ takes the value $1$ on its id\`ele.
\hfill$\square$
% END CFT_VIII613_PROOF

\section*{Original supporting excerpts}
\noindent\emph{Source \texttt{CFT.pdf}, physical page 238.}\par\medskip
% BEGIN CFT_VIII613_CONTEXT
\noindent\textsc{Theorem 1.1}
\emph{Let $K$ be a number field containing a primitive $n$th root of $1$.
A nonzero element of $K$ is an $n$th power in $K$ if it is an $n$th
power in $K_v$ for all but possibly finitely many primes $v$.}

\noindent\textsc{Proof.} Recall that, for a field $k$ containing a primitive
$n$th root of $1$, a polynomial $X^n-a$ splits completely in $k[X]$ if
it has a root $k$. Let $a$ be a nonzero element of $K$, and let $\beta$
be an $n$th root of $a$ in some extension field. If $a$ is an $n$th power
in $K_v$ for almost all $v$, then $X^n-a$ splits completely in $K_v[X]$
for almost all $v$, and so $v$ splits completely in $K[\beta]$. Since
this holds for almost all $v$, (VII, 4.6) shows that $K[\beta]=K$.
\hfill$\square$

\medskip
\noindent\emph{Source \texttt{CFT.pdf}, physical pages 261--262.}\par\medskip
\noindent\textbf{\textit{The nonarchimedean case.}} Let $K$ be a local
field. Recall that the Hilbert symbol $(\cdot,\cdot)$ can defined for
$a,b\in K^\times$ by
\[
(a,b)=\begin{cases}
1 &\Longleftrightarrow X^2-aY^2-bZ^2\text{ represents }0
  \Longleftrightarrow aY^2+bZ^2\text{ represents }1\\
-1 &\text{otherwise.}
\end{cases}
\]

\noindent\textsc{Lemma 6.5}
\emph{The Hilbert symbol has the following properties:}
\begin{enumerate}
\renewcommand{\labelenumi}{(\alph{enumi})}
\item \emph{it is bi-multiplicative and $(ac^2,bd^2)=(a,b)$ for all
$a,b,c,d\in K^\times$;}
\item \emph{for any nonsquare $a\in K^\times$, there exists a
$b\in K^\times$ such that $(a,b)=-1$;}
\item \emph{$(b,a)=(a,b)^{-1}=(a,b)$;}
\item \emph{$(a,-a)=(1,a)=1$.}
\end{enumerate}
\noindent\textsc{Proof.} Obviously, $(a,b)$ does not change when $a$ or
$b$ is multiplied by a square. Also, (c) is obvious.

Note that $(a,b)=1$ if and only if $b$ is a norm from $K[\sqrt{a}]$.
From local class field theory, we know that if $a$ is not a square in
$K$, then $\operatorname{Nm}(K[\sqrt{a}]^\times)$ is a subgroup of
index 2 in $K^\times$, and therefore $b\longmapsto(a,b)$ is an
isomorphism $K^\times/\operatorname{Nm}(K[\sqrt{a}]^\times)
\longrightarrow\{\pm1\}$. This completes the proof of (a) and (b).
Finally, $aX^2-aY^2=a(X^2-Y^2)$, and $X^2-Y^2$ represents $a^{-1}$
because it represents $0$. \hfill$\square$
% END CFT_VIII613_CONTEXT

\medskip
\noindent\emph{The fixed reading for the global reciprocity input,
Chapter VII Theorem 8.1, is retained separately as
\texttt{900008207/r001}, from physical page 230.}
\end{document}
